%%%%%%%%%%%%%%%%%%%%%%%%%%%%%%%%%%%%%%%%%%%%%%%%%%%%%%%%%%%%%%%%%%%%%%%%%%%%%%%
%% Apêndice B --- Práticas de sistemas críticos
%%
%% Absorve o relatório de práticas de sistemas críticos. Justifica escolhas de
%% engenharia de firmware; não toca a tese sobre a cadeia de medição.
%%
%% PENDENTE DE TRANSCRIÇÃO: a tabela de mecanismos por setor, com as colunas
%% Limite conhecido e Situação no Kaelix, está em
%% docs/relatorio/praticas-sistemas-criticos.tex e deve ser transportada.
%%%%%%%%%%%%%%%%%%%%%%%%%%%%%%%%%%%%%%%%%%%%%%%%%%%%%%%%%%%%%%%%%%%%%%%%%%%%%%%

\chapter{Práticas de sistemas críticos aplicadas ao firmware}
\label{ape:praticas}

Este apêndice registra a origem documentada das práticas de engenharia adotadas no
firmware. A unidade de análise é o \emph{mecanismo}, e não a norma: uma norma aparece como
origem de vários mecanismos.

Duas ressalvas delimitam o alcance. Primeira: aviônica civil e automotivo entram como
origem documentada de prática, não como alvo de certificação --- o dispositivo é protótipo
acadêmico e nenhum processo de certificação é reivindicado. Segunda: para cada mecanismo
registra-se também o que ele comprovadamente \emph{não} cobre; sem essa coluna, o
levantamento viraria lista de recomendações.

\section{Escalas de criticidade}
\label{sec:ape-escalas}

Não há equivalência normativa entre as escalas setoriais, e a correspondência citada na
literatura é informal. A norma de aviônica classifica software pela severidade da condição
de falha e não atribui taxa de falha ao software \cite{do178c}; a norma genérica é a única
das três com faixa numérica direta de medida de falha alvo \cite{iec61508}; a automotiva
compõe severidade, exposição e controlabilidade \cite{iso26262}.

\section{Mecanismos adotados}
\label{sec:ape-mecanismos}

Os mecanismos concentram-se em disciplina de codificação, detecção de erro e definição de
estado seguro --- categorias cujo custo é de projeto, e não de replicação de hardware, e
que portanto são acessíveis a um protótipo de baixo custo.

O mecanismo com consequência mais direta sobre este trabalho é a propagação explícita de
estado não implementado. Um caminho ainda não escrito devolve erro e propaga-o até o
receptor, em vez de retornar sucesso ou falhar silenciosamente. A justificativa é a mesma
que a literatura de sistemas críticos oferece para a redundância diversa: falha silenciosa
é indistinguível de operação correta, e é essa indistinguibilidade que remove a evidência
de que algo está errado \cite{knight1986,esa1996}.

\section{Limite conhecido das práticas adotadas}
\label{sec:ape-limites}

A disciplina de codificação e a propagação de erro não substituem redundância de hardware,
e a literatura documenta que a suposição de independência entre versões diversas de
software não se sustenta empiricamente \cite{knight1986}. \emandamento{Nenhuma das práticas
adotadas foi exercitada em hardware real; sua eficácia é verificada apenas em teste no
\emph{host}.}
